% Copyright (c) 2026 Sergey Zhunkov.
% SPDX-License-Identifier: CC-BY-4.0
%
% Set \ChartLanguage to en or ru before loading this source. English is the default.
\documentclass[a4paper,landscape]{article}
\usepackage[margin=0pt]{geometry}
\usepackage{fontspec,babel,etoolbox,tikz,luacode,hyperref,bookmark}
\usetikzlibrary{babel}

\providecommand{\ChartLanguage}{en}
\newif\ifChartRussian
\ifdefstring{\ChartLanguage}{ru}{%
  \ChartRussiantrue
  \babelprovide[import,main]{russian}
  \babelprovide[import]{english}
}{%
  \ifdefstring{\ChartLanguage}{en}{%
    \babelprovide[import,main]{english}
    \babelprovide[import]{russian}
  }{%
    \PackageError{capo}{Unsupported ChartLanguage '\ChartLanguage'}%
      {Set ChartLanguage to en or ru.}%
  }%
}
\newcommand{\Localize}[2]{\ifChartRussian#1\else#2\fi}

\setmainfont{DejaVuSansCondensed.ttf}[BoldFont=DejaVuSansCondensed-Bold.ttf]
\pagestyle{empty}
\definecolor{ChartInk}{RGB}{23,43,50}
\definecolor{ChartMuted}{RGB}{70,91,99}
\definecolor{ChartRule}{RGB}{197,208,209}
\definecolor{ChartStripe}{RGB}{244,247,247}
\providecommand{\ChartAccentHex}{195755}
\definecolor{ChartAccent}{HTML}{\ChartAccentHex}
\colorlet{ChartHeader}{ChartAccent!13!white}
\colorlet{ChartHeading}{ChartAccent!33!ChartInk}

\newcommand{\ChartTitle}{\Localize{Каподастр на гитаре}{Guitar capo}}
\hypersetup{pdfauthor={Sergey Zhunkov},pdftitle={\ChartTitle},
  colorlinks=true,urlcolor=ChartHeading,linkcolor=ChartHeading}
\ifChartRussian
  \hypersetup{pdflang={ru-RU}}
\else
  \hypersetup{pdflang={en}}
\fi

\begin{luacode*}
-- Root pitch classes start at C = 0. The suffix travels with the shape.
local roots = {
  [""]  = {"C", "D♭", "D", "E♭", "E", "F", "F♯", "G", "A♭", "A", "B♭", "B"},
  ["m"] = {"C", "C♯", "D", "E♭", "E", "F", "F♯", "G", "G♯", "A", "B♭", "B"},
}
local shapes = {
  {0, ""}, {9, ""}, {7, ""}, {4, ""}, {2, ""},
  {9, "m"}, {2, "m"}, {4, "m"},
}

function PrintCapoLookup()
  for column, shape in ipairs(shapes) do
    local root, suffix = table.unpack(shape)
    local names = roots[suffix]
    tex.sprint(string.format([[\LookupShape{%d}{%s}]], column, names[root + 1] .. suffix))
    for fret = 0, 12 do
      local sounding = names[(root + fret) % 12 + 1] .. suffix
      tex.sprint(string.format([[\LookupCell{%d}{%d}{%s}]], fret, column, sounding))
    end
  end
end
\end{luacode*}

% Coordinates and font sizes are in PDF points, measured from the bottom left.
\newcommand{\ChartText}[6][base west]{%
  \node[anchor=#1,inner sep=0pt,outer sep=0pt,text=#5,
    font=\fontsize{#4bp}{#4bp}\selectfont] at ({#2},{#3}) {#6};%
}
\newcommand{\ChartParagraph}[4]{%
  \node[anchor=north west,inner sep=0pt,outer sep=0pt,text=ChartInk,align=left,
    text width=#3bp,font=\fontsize{11bp}{14bp}\selectfont] at ({#1},{#2}) {#4};%
}
\newcommand{\ChartRuleAt}[1]{\draw[ChartRule,line width=.65bp] (34,#1) -- (808,#1);}
\newcommand{\ChartHeadingAt}[2]{\ChartText{34}{#1}{15}{ChartHeading}{\bfseries #2}}
\newcommand{\Card}[4]{%
  \draw[ChartRule,line width=.65bp,rounded corners=2bp] (#1,#2) rectangle (#3,#4);%
}
\newcommand{\ReferenceLink}[2]{%
  \href{https://github.com/zhunkoff/music-reference/blob/main/\Localize{README.ru.md}{README.md}\#ref-#1}{#2}%
}
\newcommand{\ChartPage}[3]{%
  \pdfbookmark[0]{#2}{capo-\thepage}%
  \AddToHookNext{shipout/background}{%
    \put(0,-\paperheight){\begin{tikzpicture}[overlay,x=1bp,y=1bp]
      \ChartText{34}{542}{28}{ChartInk}{\bfseries\ChartTitle}%
      \ChartText[base east]{808}{542}{12}{ChartHeading}{\bfseries #1 / 2 · #2}%
      \ChartRuleAt{529}%
      #3%
      \ChartRuleAt{43}%
      \ChartText{34}{27}{9.5}{ChartMuted}{%
        \ReferenceLink{guitar-chords}{\Localize{Гитарные аккорды}{Guitar chords}} ·
        \ReferenceLink{transposition}{\Localize{Транспонирование}{Transposition}}\qquad
        \textbf{\Localize{Источник}{Source}}\quad
        \href{https://www.daddario.com/globalassets/pdfs/accessories/daddario_capo_guide.pdf}
          {Marc Schonbrun · The D’Addario Capo Guide}}%
      \ChartText[base east]{808}{27}{9}{ChartMuted}{%
        © 2026 \href{https://github.com/zhunkoff/music-reference}{Sergey Zhunkov} ·
        \href{https://creativecommons.org/licenses/by/4.0/}{CC BY 4.0}}%
    \end{tikzpicture}}%
  }%
  \null\newpage
}

\newcommand{\LookupX}[1]{146+82.75*(#1-.5)}
\newcommand{\LookupY}[1]{413-19*(#1+.5)}
\newcommand{\LookupShape}[2]{\ChartText[base]{\LookupX{#1}}{420}{12}{ChartHeading}{\bfseries #2}}
\newcommand{\LookupCell}[3]{\ChartText[base]{\LookupX{#2}}{\LookupY{#1}-4}{12}{ChartInk}{#3}}
\newcommand{\CapoLookup}{%
  \fill[ChartHeader] (34,413) rectangle (808,434);%
  \foreach \fret in {0,...,12}{%
    \ifodd\fret\relax
      \fill[ChartStripe] (34,{\LookupY{\fret}-9.5}) rectangle (808,{\LookupY{\fret}+9.5});%
    \fi
    \ChartText[base]{90}{\LookupY{\fret}-4}{11.5}{ChartHeading}{\bfseries\fret}%
  }%
  \draw[ChartRule,line width=.65bp] (146,166) -- (146,434) (34,413) -- (808,413);%
  \Card{34}{166}{808}{434}%
  \ChartText[base]{90}{420}{11}{ChartHeading}{\bfseries\Localize{Лад каподастра}{Capo fret}}%
  \directlua{PrintCapoLookup()}%
}

\newcommand{\ProgressionRow}[7]{%
  \ChartText{48}{#1}{11}{ChartHeading}{\bfseries #2}%
  \ChartText[base]{380}{#1}{12}{ChartInk}{#3}%
  \ChartText[base]{475}{#1}{12}{ChartInk}{#4}%
  \ChartText[base]{570}{#1}{12}{ChartInk}{#5}%
  \ChartText[base]{665}{#1}{12}{ChartInk}{#6}%
  \ChartText[base]{760}{#1}{12}{ChartInk}{#7}%
}

\begin{document}
\ChartPage{1}{\Localize{Таблица звучащих аккордов}{Sounding-chord lookup}}{%
  \ChartParagraph{34}{510}{774}{\Localize
    {Полный каподастр прижимает все шесть струн. Стандартный строй (струны 6→1) — E2–A2–D3–G3–B3–E4.
     Сдвиг каподастра на лад вверх повышает все звуки формы на полутон. Лад каподастра считают
     от верхнего порожка, лады формы — от каподастра.}
    {A full capo clamps all six strings. Standard tuning without it, from string 6 to 1, is
     E2–A2–D3–G3–B3–E4. Moving the capo up one fret raises every pitch of the same shape by one semitone.
     Count the physical capo fret from the nut. In a fingering, the capo becomes fret zero.}}%

  \ChartHeadingAt{460}{\Localize{Звучащий аккорд по ладу каподастра и форме}{Sounding chord by capo fret and shape}}%
  \ChartParagraph{34}{453}{774}{\Localize
    {Найдите форму в заголовке и физический лад в левом столбце. Строка 0 означает игру без каподастра.}
    {Find the shape across the top and the physical capo fret on the left. Row 0 means no capo.}}%
  \CapoLookup
  \ChartParagraph{34}{155}{774}{\Localize
    {Строка 12 повторяет названия из строки 0. Каждый звук на октаву выше.
     Названия международные (B = си). В парах D♭/C♯, E♭/D♯, F♯/G♭, A♭/G♯ и B♭/A♯
     оба обозначения соответствуют одному классу высоты. Выбирайте написание по тональности.}
    {At fret 12 the chord names repeat, with every pitch an octave above row 0. Names use the international
     C–B system. D♭/C♯, E♭/D♯, F♯/G♭, A♭/G♯, and B♭/A♯ are enharmonic pairs. Choose the spelling that fits the key.}}%

  \Card{34}{60}{808}{114}%
  \ChartParagraph{48}{103}{746}{\Localize
    {\textbf{Форма C, каподастр на 2-м ладу.} Лады формы от 6-й струны к 1-й записаны как x–3–2–0–1–0.
     Здесь x — приглушённая струна, 0 — струна, которую прижимает только каподастр. На 5-й струне лад 3 от каподастра
     соответствует физическому ладу 5. Звуки C3–E3–G3–C4–E4 повышаются до D3–F♯3–A3–D4–F♯4 (мажорное трезвучие D).}
    {\textbf{C shape, capo at fret 2.} The shape’s frets from string 6 to 1 are x–3–2–0–1–0,
     where x is muted and 0 is open at the capo. On string 5, relative fret 3 is physical fret 5.
     C3–E3–G3–C4–E4 rises to D3–F♯3–A3–D4–F♯4 (a D major chord).}}%
}

\ChartPage{2}{\Localize{Подбор форм и последовательности}{Reverse lookup and progressions}}{%
  \ChartHeadingAt{505}{\Localize{От звучащей тональности к формам}{From a sounding key to shapes}}%
  \ChartParagraph{34}{493}{774}{\Localize
    {Выпишите звучащие аккорды целевой тональности или конкретной последовательности.
     Начните с тоники или другого нужного аккорда. Найдите его в таблице с учётом качества
     (мажор или минор) и энгармонии. Для каждого совпадения запишите лад и форму. В той же строке
     проверьте остальные аккорды. Выберите лад, на котором все нужные формы удобно сыграть.}
    {List the sounding chords required in the target key or progression. Start with its tonic or another required chord.
     Find that chord in the table, matching its quality (major or minor) and allowing for enharmonic names.
     For each match, note the fret and shape. Check the remaining chords in that same row.
     Choose a fret where all the required shapes are playable for you.}}%
  \ChartParagraph{34}{430}{774}{\Localize
    {\textbf{Один аккорд.} Для звучащего E♭ подходят, например, форма D на 1-м ладу и форма C на 3-м.
     Если цель — ми-бемоль мажор, проверьте в каждой из этих строк и другие аккорды песни.}
    {\textbf{One chord.} For a sounding E♭ major chord, two options are a D shape at fret 1 and a C shape at fret 3.
     If the target key is E♭ major, check the other chords in the song in each candidate row.}}%

  \Card{34}{253}{808}{393}%
  \ChartText{48}{374}{13}{ChartHeading}{\bfseries\Localize{Ля мажор · каподастр на 2-м ладу}{A major · capo at fret 2}}%
  \fill[ChartHeader] (35,343) rectangle (807,361);%
  \ProgressionRow{348}{\Localize{Ступени}{Scale degrees}}{I}{IV}{V}{vi}{I}%
  \ProgressionRow{329}{\Localize{Формы}{Shapes}}{G}{C}{D}{Em}{G}%
  \ProgressionRow{310}{\Localize{Звучащие аккорды}{Sounding chords}}{A}{D}{E}{F♯m}{A}%
  \ChartParagraph{48}{294}{746}{\Localize
    {На 5-м ладу форма E тоже даёт тонику A. В этой строке нет F♯m. На 2-м ладу есть все четыре
     нужных аккорда. Форма Em даёт F♯m, сохраняя минорное качество.}
    {An E shape at fret 5 also gives the tonic A. That row has no F♯m. Fret 2 covers all four required chords.
     The Em shape gives F♯m and keeps its minor quality.}}%

  \Card{34}{64}{808}{237}%
  \ChartText{48}{218}{13}{ChartHeading}{\bfseries\Localize
    {Из ля минора в соль минор · каподастр на 3-м ладу}
    {From A minor to G minor · capo at fret 3}}%
  \fill[ChartHeader] (35,187) rectangle (807,205);%
  \ProgressionRow{192}{\Localize{Ступени натурального минора}{Natural minor scale degrees}}{i}{iv}{VII}{III}{i}%
  \ProgressionRow{173}{\Localize{Ля минор без каподастра}{A minor without capo}}{Am}{Dm}{G}{C}{Am}%
  \ProgressionRow{154}{\Localize{Формы с каподастром на 3-м ладу}{Shapes with capo at fret 3}}{Em}{Am}{D}{G}{Em}%
  \ProgressionRow{135}{\Localize{Звучащие аккорды в соль миноре}{Sounding chords in G minor}}{Gm}{Cm}{F}{B♭}{Gm}%
  \ChartParagraph{48}{119}{746}{\Localize
    {Целевая тональность на два полутона ниже исходной. Возьмите формы последовательности в ми миноре
     и проверьте их по строке 3. Каподастр поднимает эти формы на три полутона до соль минора.
     Для перехода в более низкую тональность здесь меняются формы. Регистр отдельных звуков может измениться.}
    {The target key is two semitones below the original. Use the shapes of the progression in E minor and check them
     against row 3. The capo raises those shapes by three semitones into G minor. The change of shapes makes the lower
     sounding key possible. Individual notes may sit in different octaves.}}%
}
\end{document}
